% rhythm: article text blocks: paragraph to paragraph, paragraph to and from both heading levels, heading to heading, and where the first and last lines stand on the page
% boundary: page-first-baseline Alder top
% boundary: par-par Birch span
% boundary: par-section Cedar span
% boundary: section-par Dogwood span
% boundary: par-par Elm span
% boundary: par-subsection Hazel span
% boundary: subsection-par Juniper span
% boundary: par-section Larch span
% boundary: section-subsection Maple span
% boundary: subsection-par Oak span
% boundary: page-footline =1 top
\documentclass{article}
\usepackage{fontspec}
\setmainfont{OpenSans-Regular.ttf}[Path=../corpus/fonts/, BoldFont=OpenSans-Bold.ttf, ItalicFont=OpenSans-Italic.ttf, BoldItalicFont=OpenSans-BoldItalic.ttf]
\begin{document}
Alder opens the page with one short line.

Birch follows as a second paragraph of one line.

\section{Cedar heading}
Dogwood starts the paragraph under the heading and runs on long enough to
wrap onto a second line, so that the boundary below it is read from the
last line of a wrapped paragraph rather than from its first line.

Elm is a plain paragraph between two others.

\subsection{Hazel subheading}
Juniper sits under the subsection heading.

\section{Larch heading}
\subsection{Maple subheading}
Oak closes the page.
\end{document}
